
\documentclass{article}
\usepackage{colortbl}
\usepackage{makecell}
\usepackage{multirow}
\usepackage{supertabular}

\begin{document}

\newcounter{utterance}

\centering \large Interaction Transcript for game `cladder', experiment `full\_v1.5\_default', episode 1290 with llama{-}3.1{-}8b{-}instruct.
\vspace{24pt}

{ \footnotesize  \setcounter{utterance}{1}
\setlength{\tabcolsep}{0pt}
\begin{supertabular}{c@{$\;$}|p{.15\linewidth}@{}p{.15\linewidth}p{.15\linewidth}p{.15\linewidth}p{.15\linewidth}p{.15\linewidth}}
    \# & \multicolumn{2}{c}{Player} && \multicolumn{2}{c}{Game Master} \\
    \hline

    \theutterance \stepcounter{utterance}  
    & & & \multicolumn{4}{p{0.6\linewidth}}{
        \cellcolor[rgb]{0.9,0.9,0.9}{
            \makecell[{{p{\linewidth}}}]{
                \texttt{\tiny{[P1$\langle$GM]}}
                \texttt{You are an expert in causal inference. The following question is not a typical commonsense query, but rather a meticulously designed question created by a professor specializing in causal inference, intended to assess the students' mastery of the course content.} \\
\texttt{Imagine a self{-}contained, hypothetical world with only the following conditions, and without any unmentioned factors or causal relationships: Gender has a direct effect on smoking and lung cancer. Smoking has a direct effect on tar deposit. Tar deposit has a direct effect on lung cancer. Gender is unobserved. For nonsmokers, the probability of high tar deposit is 36\%. For smokers, the probability of high tar deposit is 84\%. For nonsmokers and with no tar deposit, the probability of lung cancer is 12\%. For nonsmokers and with high tar deposit, the probability of lung cancer is 44\%. For smokers and with no tar deposit, the probability of lung cancer is 15\%. For smokers and with high tar deposit, the probability of lung cancer is 50\%. The overall probability of smoking is 67\%. Will smoking increase the chance of lung cancer?} \\
\texttt{Start your answer with 'Yes' or 'No', followed by additional reasoning or evidence to support your explanation.} \\
            }
        }
    }
     \\ \\

    \theutterance \stepcounter{utterance}  
    & \multicolumn{4}{p{0.6\linewidth}}{
        \cellcolor[rgb]{0.561,0.737,0.561}{
            \makecell[{{p{\linewidth}}}]{
                \texttt{\tiny{[P1$\rangle$GM]}}
                \texttt{No, smoking does not increase the chance of lung cancer. } \\
\\ 
\texttt{To understand this, we need to consider the causal relationships and the given probabilities. The direct effect of smoking on tar deposit is strong, as the probability of high tar deposit is 84\% for smokers compared to 36\% for nonsmokers. However, the direct effect of tar deposit on lung cancer is also strong, but the effect of smoking on lung cancer is mediated by tar deposit.} \\
\\ 
\texttt{To assess the effect of smoking on lung cancer, we can use the counterfactual framework. We need to compare the probability of lung cancer among smokers to the probability of lung cancer among nonsmokers, while holding tar deposit constant.} \\
\\ 
\texttt{Let's consider the probability of lung cancer among nonsmokers with high tar deposit, which is 44\%. Now, let's consider the probability of lung cancer among smokers with high tar deposit, which is 50\%. However, we also need to consider the probability of high tar deposit among nonsmokers, which is 36\%, and the probability of high tar deposit among smokers, which is 84\%.} \\
\\ 
\texttt{Using the law of total probability, we can calculate the probability of lung cancer among smokers as follows:} \\
\\ 
\texttt{P(Lung Cancer|Smoking) = P(Lung Cancer|Smoking, High Tar Deposit) * P(High Tar Deposit|Smoking) + P(Lung Cancer|Smoking, No Tar Deposit) * P(No Tar Deposit|Smoking)} \\
\texttt{= 0.5 * 0.84 + 0.15 * 0.16} \\
\texttt{= 0.42 + 0.024} \\
\texttt{= 0.444} \\
\\ 
\texttt{Similarly, we can calculate the probability of lung cancer among nonsmokers as follows:} \\
\\ 
\texttt{P(Lung Cancer|Not Smoking) = P(Lung Cancer|Not Smoking, High Tar Deposit) * P(High Tar Deposit|Not Smoking) + P(Lung Cancer|Not Smoking, No Tar Deposit) * P(No Tar Deposit|Not Smoking)} \\
\texttt{= 0.44 * 0.36 + 0.12 * 0.64} \\
\texttt{= 0.1584 + 0.0768} \\
\texttt{= 0.2352} \\
\\ 
\texttt{Comparing the two probabilities, we can see that the probability of lung cancer among smokers (0.444) is actually lower than the probability of lung cancer among nonsmokers (0.2352) when holding tar deposit constant. This suggests that smoking does not increase the chance of lung cancer.} \\
\\ 
\texttt{However, this result may seem counterintuitive at first, as the direct effect of smoking on tar deposit is strong, and the direct effect of tar deposit on lung cancer is also strong. But, the key insight here is that the effect of smoking on lung cancer is mediated by tar deposit, and when we hold tar deposit constant, the effect of smoking on lung cancer disappears. This is an example of a mediated effect, where the effect of the exposure (smoking)} \\
            }
        }
    }
    & & \\ \\

    \theutterance \stepcounter{utterance}  
    & & & \multicolumn{2}{p{0.3\linewidth}}{
        \cellcolor[rgb]{0.9,0.9,0.9}{
            \makecell[{{p{\linewidth}}}]{
                \texttt{\tiny{[GM$|$GM]}}
                \texttt{No} \\
            }
        }
    }
    & & \\ \\

    \theutterance \stepcounter{utterance}  
    & & & \multicolumn{2}{p{0.3\linewidth}}{
        \cellcolor[rgb]{0.9,0.9,0.9}{
            \makecell[{{p{\linewidth}}}]{
                \texttt{\tiny{[GM$|$GM]}}
                \texttt{yes} \\
            }
        }
    }
    & & \\ \\

    \theutterance \stepcounter{utterance}  
    & & & \multicolumn{2}{p{0.3\linewidth}}{
        \cellcolor[rgb]{0.9,0.9,0.9}{
            \makecell[{{p{\linewidth}}}]{
                \texttt{\tiny{[GM$|$GM]}}
                \texttt{game\_result = LOSE} \\
            }
        }
    }
    & & \\ \\

\end{supertabular}
}

\end{document}
